\subsection{Sprint 2}

Após a retrospectiva da Sprint 1 e o grande sucesso que ela foi, o time estava bastante entusiasmado com o projeto em geral, preparamos com cautela nossos próximos passos uma vez que a Sprint que estaria por vir teria muitas features chaves em nosso projeto. Foi algo muito agradável ver meus colegas, em geral, contentes com o trabalho que produziram ao longo do projeto, tendo em vista que grande parte do time era AGES I. Fiquei impressionado com a produtividade em geral do time mostrando
que aparentemente todos que estavam lá estavam realmente dispostos a contribuir com o projeto.

Para esta Sprint, precisaríamos implementar operações relacionadas ao cadastro de certificados, o que exigiria que nossa operação na blockchain estivesse funcionando perfeitamente. A lógica da aplicação seria a seguinte: o time da blockchain decidiu trabalhar com uma sidechain do Ethereum. Teríamos uma API da blockchain rodando em Solidity (linguagem de programação do Ethereum), que seria responsável por realizar as diversas operações na blockchain.

Para a minha equipe, a do backend, nosso trabalho seria principalmente fazer a "conexão" entre o frontend e a blockchain. Para isso, precisaríamos desenvolver uma API RESTful que permitisse a comunicação eficiente e segura entre as duas camadas. Isso envolveria a criação de endpoints para registro, consulta e validação de certificados, além da implementação de mecanismos de autenticação e autorização. Além das implementações acima, precisaríamos também implementar nosso banco de dados no qual eu e meu colega AGES II estavamos quase terminando sua modelagem.

No início da Sprint, tudo correu como esperado. Conseguimos concluir a modelagem do banco de dados por completo, o que nos permitiu começar a implementar a lógica do banco de dados (relacionamentos, atributos) no backend. Apesar do progresso que estávamos fazendo, as enchentes que ocorreram no Rio Grande do Sul mudaram completamente o rumo das coisas. Ficamos duas semanas sem aulas e, durante esse período, pausamos o projeto para lidar com os estragos deixados pela enchente, tanto materiais quanto psicológicos.

Após essas duas semanas sem aulas, elas retornaram de forma online. Nossa primeira reunião pós-enchente foi um pouco melancólica, pois era evidente que a moral do time, em geral, estava muito baixa, inclusive a minha. Embora eu não tenha sofrido danos materiais durante a enchente, os danos psicológicos foram aparentes para mim, o que levou a uma piora no meu quadro de doença autoimune.

Devido à minha baixa moral e à piora na minha saúde, não consegui contribuir para o projeto até o final da sprint. Foi algo bastante triste, pois eu estava muito comprometido e animado com a AGES. Ter que parar e perder o ritmo do projeto e não conseguir acompanhar meus colegas foi desanimador e impactou negativamente minha motivação até o final do projeto.

Apesar da minha ausência e de todos os problemas que ocorreram nesta sprint, o time impressionantemente foi capaz de fazer bastante, fico muito orgulhoso do time e com o comprometimento que eles tiveram para terminar a sprint da melhor forma possível, mesmo que não fosse da mesma maneira que foi a passada.

A sprint teve seu fim no dia 24/05/2024 com uma videochamada que fizemos com o stakeholder. Apesar de alguns débitos técnicos em relação às funcionalidades do backend, à configuração da AWS e à camada de conexão da blockchain ainda incompleta, Edimar, nosso stakeholder, ficou satisfeito com o que produzimos durante a sprint. Ele foi muito compreensivo em relação aos débitos técnicos e entendeu a situação pela qual todos nós estávamos passando, o que foi muito legal da parte dele. O time, apesar de tudo, ficou contente com o resultado dessa sprint e depois da reunião com o Edimar ficou perceptível que meus colegas começaram a recuperar aquele ânimo perdido após as enchentes.

A sprint, apesar de todos os percalços que tivemos durante sua execução, nos ensinou muitas coisas. Embora tenha sido a pior em termos de entregas, foi a que mais nos ensinou lições valiosas sobre resiliência, adaptação e trabalho em equipe. Aprendemos a importância de apoiar uns aos outros em tempos difíceis, a necessidade de ajustar nossos planos diante de imprevistos e a capacidade de manter o foco no objetivo final, mesmo quando as circunstâncias não são ideais.
